\appendix

\section{Perturbative Partition, Channel Decomposition, and Contact Transformation}

In the main text the theory is organized in terms of blocks $H_{nm}$, where
$n$ and $m$ denote the vibrational and rotational degree, respectively. It is
important to distinguish this bidegree from the perturbative order assigned to
each block. The bidegree is a structural label, whereas the perturbative order
depends on the adopted Hamiltonian partition.

Throughout the present work the standard Watson/VPT2-style partition is
assumed. Within this convention,
\begin{equation}
H_{03},\;H_{12},\;H_{30}
\end{equation}
are first-order terms, whereas
\begin{equation}
H_{40},\;H_{22},\;H_{13},\;H_{31},\;H_{04}
\end{equation}
are second-order terms. This convention is standard in Watson-style
rovibrational perturbation theory, but it is not universal. In alternative
state-summation formulations, which are also often labeled as VPT2, the same
blocks may be counted differently.

Let the transformed Hamiltonian be written in Van Vleck form as
\begin{equation}
\widetilde H = e^{iS} H e^{-iS},
\end{equation}
with $S=S^{(1)}+S^{(2)}+\cdots$. The block-resolved effective Hamiltonian is
then obtained by expanding the Baker--Campbell--Hausdorff series,
\begin{equation}
\widetilde H
=
H + i[S,H] + \frac{i^2}{2!}[S,[S,H]] + \frac{i^3}{3!}[S,[S,[S,H]]] + \cdots .
\end{equation}
The channel decomposition used in the main text is exact only after fixing
\begin{enumerate}
\item the perturbative partition,
\item the truncation order,
\item the specific Van Vleck/contact-transformation construction.
\end{enumerate}
The coefficients of the channels are therefore not determined by the block
labels alone, but also by the combinatorial prefactors and factorial
denominators generated by the BCH series.

At the level relevant for quartic effective constants, the standard channels
may be organized schematically as
\begin{equation}
H_{12}H_{12},\qquad H_{22},\qquad H_{12}H_{30},\qquad H_{30}H_{30},\qquad H_{04},
\end{equation}
with the understanding that each label denotes the fully weighted contribution
generated by the adopted perturbative construction, not a bare symbolic
product.

\section{Quartic Tensor Reconstruction and Pseudoinverse Gauge Fixing}

\subsection{Reduced quartic tensor}

The quartic centrifugal distortion sector is most naturally described in terms
of the pair--pair quartic tensor
\begin{equation}
\boldsymbol\tau^{(4)} =
\bigl(
\tau_{aaaa},\tau_{bbbb},\tau_{cccc},
\tau_{aabb},\tau_{aacc},\tau_{bbcc}
\bigr)^T ,
\end{equation}
or equivalently by the reduced symmetric matrix
\begin{equation}
\tau'=
\begin{pmatrix}
\tau_{aaaa} & \tau_{aabb}/2 & \tau_{aacc}/2 \\
\tau_{aabb}/2 & \tau_{bbbb} & \tau_{bbcc}/2 \\
\tau_{aacc}/2 & \tau_{bbcc}/2 & \tau_{cccc}
\end{pmatrix}.
\end{equation}
This tensor contains six independent real components.

In the standard harmonic problem, the corresponding quartic tensor is generated
by the first derivative of the inverse inertia tensor with respect to the
normal coordinates,
\begin{equation}
\mu_{1,k} \equiv \frac{\partial I^{-1}}{\partial Q_k},
\end{equation}
through the Wilson quartic tensor
\begin{equation}
\tau_{\alpha\beta\gamma\delta}
=
-\frac12
\sum_k
\frac{
\mu_{\alpha\beta}^{(k)}
\mu_{\gamma\delta}^{(k)}
}{\omega_k^2}.
\label{eq:app_quartic_tau}
\end{equation}
In the implementation used in the accompanying code base, this same quartic
content may also be written in Coriolis form. The two formulations are
equivalent in the standard quartic sector.

\subsection{Forward projection to Watson quartic constants}

Once $\tau'$ has been reconstructed, one first defines
\begin{equation}
T_{ij} = \frac{\tau'_{ij}}{4},
\end{equation}
then the cylindrical combinations
\begin{align}
T_{400} &= \frac{3T_{aa}+3T_{bb}+2T_{ab}}{8}, \\
T_{220} &= T_{ac}+T_{bc}-2T_{400}, \\
T_{040} &= T_{cc}-T_{220}-T_{400}, \\
T_{202} &= \frac{T_{aa}-T_{bb}}{4}, \\
T_{022} &= \frac{T_{ac}-T_{bc}}{2}-T_{202}, \\
T_{004} &= \frac{T_{aa}+T_{bb}-2T_{ab}}{16}.
\end{align}
For the $A$ reduction, the quartic constants are
\begin{align}
\Delta_J   &= -T_{400}-2T_{004}, \\
\Delta_{JK}&= -T_{220}+12T_{004}, \\
\Delta_K   &= -T_{040}-10T_{004}, \\
\delta_J   &= -T_{202}, \\
\delta_K   &= -T_{022}-4\Sigma T_{004},
\end{align}
with
\begin{equation}
\Sigma = \frac{2A-B-C}{B-C}.
\end{equation}
For the $S$ reduction one has
\begin{align}
D_J   &= -T_{400}+\frac12 T_{022}\Sigma_1, \\
D_{JK}&= -T_{220}-3T_{022}\Sigma_1, \\
D_K   &= -T_{040}+\frac52 T_{022}\Sigma_1, \\
d_1   &= T_{202}, \\
d_2   &= T_{004}+\frac14 T_{022}\Sigma_1,
\end{align}
with
\begin{equation}
\Sigma_1 = \frac{B-C}{2A-B-C} = \Sigma^{-1}.
\end{equation}

\subsection{Rank deficiency and affine gauge family}

The forward map from the six-dimensional pair--pair quartic tensor to the five
Watson quartic constants is rank-deficient. A fixed quartic spectroscopic
parameter set therefore determines an affine family of compatible tensors,
rather than a unique tensor representative.

In the $A$ reduction, a null vector of the quartic projection map may be
written as
\begin{equation}
\mathbf n =
\left(
0,\,
0,\,
0,\,
1,\,
\frac{\Sigma-1}{2},\,
-\frac{\Sigma+1}{2}
\right)^T.
\end{equation}
In reduced-matrix form this corresponds to the null tensor
\begin{equation}
N'=
\begin{pmatrix}
0 & 1/2 & (\Sigma-1)/4 \\
1/2 & 0 & -(\Sigma+1)/4 \\
(\Sigma-1)/4 & -(\Sigma+1)/4 & 0
\end{pmatrix}.
\end{equation}
Hence, if $\tau'$ reproduces a given quartic parameter set, so does
\begin{equation}
\tau' + \eta N'
\end{equation}
for arbitrary scalar $\eta$.

The Moore--Penrose pseudoinverse is used in the present work to fix this affine
gauge freedom. The corresponding pseudoinverse gauge condition is
\begin{equation}
2\tau'_{ab}+(\Sigma-1)\tau'_{ac}-(\Sigma+1)\tau'_{bc}=0.
\label{eq:app_quartic_gauge}
\end{equation}
Equation~\eqref{eq:app_quartic_gauge} is not a physical identity. It is the
gauge slice selected by the pseudoinverse reconstruction algorithm.

\section{The Standard Sextic Sector Without Primitive \texorpdfstring{$d^2I/dQ^2$}{d2I/dQ2}}

\subsection{Operational sextic structure}

The standard sextic centrifugal distortion constants are written in terms of
the ten Cartesian sextic coefficients
\begin{equation}
\Phi_{aaa},\Phi_{aab},\Phi_{aac},\Phi_{abb},\Phi_{abc},
\Phi_{acc},\Phi_{bbb},\Phi_{bbc},\Phi_{bcc},\Phi_{ccc},
\end{equation}
which are projected onto the usual seven Watson sextic constants in the chosen
reduction. In the implementation used here, the corresponding geometry sector
is built from the standard quantities
\begin{equation}
\tau,\qquad c_1,\qquad c_2,\qquad \zeta,
\end{equation}
together with the explicit cubic anharmonic contribution $\phi_3$.

The crucial structural point established in the code analysis is that, in the
standard sextic sector, these objects do not require second derivatives of the
inertia tensor as primitive ingredients. More precisely:
\begin{enumerate}
\item the standard quartic tensor $\tau$ is generated by the first-level
inverse-inertia derivatives $\mu_1$;
\item the standard object $c_1$ is modewise proportional to $\mu_1$;
\item the unsummed object behind $c_2$ is a Coriolis-coupled image of $\mu_1$;
\item the full currently exposed standard sextic geometry sector is regenerated
from $\mu_1$ and Coriolis couplings alone.
\end{enumerate}

\subsection{Standard sextic geometry as first-level closure}

At the level relevant for the present work, the standard sextic tensorial
content may therefore be summarized schematically as
\begin{equation}
\Phi^{(6)}_{\mathrm{std}}
=
\Phi^{(6)}_{\mathrm{geom}}(\mu_1,\zeta,\omega,\mathrm{rot})
+
\Phi^{(6)}_{\mathrm{cubic}}(\phi_3).
\label{eq:app_sextic_std_split}
\end{equation}
Equation~\eqref{eq:app_sextic_std_split} expresses the main structural result:
the standard sextic geometry sector remains in the same first-level tensorial
language as the standard quartic problem. In particular, the appearance of
second inertia derivatives in some traditional derivations does not imply that
they define an independent primitive tensorial level in the standard sextic
sector.

In the explicit working formulation used in the code, the same split is
written more concretely in terms of the quantities $c_1$, $c_2$, and $\tau$:
\begin{equation}
\Phi^{(6)}_{\mathrm{std}}
=
\Phi^{(6)}_{\mathrm{geom}}(c_1,c_2,\tau)
+
\Phi^{(6)}_{\mathrm{cubic}}(\phi_3).
\label{eq:app_sextic_std_split_c1c2}
\end{equation}
The function $\Phi^{(6)}_{\mathrm{geom}}$ contains no primitive second
derivative of the inertia tensor. This is the precise sense in which the
standard sextic geometry sector closes on the same first-level tensorial
language as the standard quartic problem.

\subsection{Splitting the cubic-force contribution}

For the present paper it is useful to refine the explicit cubic term further.
Let the reduced cubic force field in normal coordinates be written as
\begin{equation}
\phi_{ijk}
=
\phi^{\mathrm{sd}}_{ijk}
+
\phi^{(3)}_{ijk},
\label{eq:app_phi3_split}
\end{equation}
where $\phi^{\mathrm{sd}}_{ijk}$ denotes the semi-diagonal part, i.e.\ the
part with at least two equal mode indices, and $\phi^{(3)}_{ijk}$ denotes the
genuinely three-index remainder, i.e.\ the part with three distinct mode
indices.

With this notation the cubic sextic contribution splits as
\begin{equation}
\Phi^{(6)}_{\mathrm{cubic}}
=
\Phi^{(6)}_{\mathrm{cubic,sd}}
+
\Phi^{(6)}_{\mathrm{cubic,3ind}},
\label{eq:app_sextic_cubic_split}
\end{equation}
with
\begin{align}
\Phi^{(6)}_{\mathrm{cubic,sd}}
&=
\Phi^{(6)}_{\mathrm{cubic}}(\phi^{\mathrm{sd}}),
\\
\Phi^{(6)}_{\mathrm{cubic,3ind}}
&=
\Phi^{(6)}_{\mathrm{cubic}}(\phi^{(3)}).
\end{align}

At the level of the implemented formulas, the cubic sector is linear in the
reduced cubic force constants, and therefore this split is exact. If
\begin{equation}
\Phi^{(6)}_{\mathrm{cubic}}
\sim
\sum_{ijk}\phi_{ijk}\,T_{ijk}(c_1),
\end{equation}
where $T_{ijk}(c_1)$ denotes the appropriate contraction of the cubic force
constants with the first-level tensors $c_1$, then
\begin{equation}
\sum_{ijk}\phi_{ijk}\,T_{ijk}
=
\sum_{ijk}\phi^{\mathrm{sd}}_{ijk}\,T_{ijk}
+
\sum_{ijk}\phi^{(3)}_{ijk}\,T_{ijk}.
\label{eq:app_sextic_cubic_linear_split}
\end{equation}

This decomposition is useful for two reasons. First, it makes explicit that
the standard sextic problem is not only ``geometry plus cubic'', but
``geometry plus semi-diagonal cubic plus genuinely three-index cubic
remainder''. Second, it adapts the sextic formulas to the computational
hierarchy relevant for medium-sized molecules, for which these two cubic-force
sectors may be available at different levels of theory.

This computational distinction is especially transparent when the cubic force
field is obtained by finite differences of analytic gradients. In that case,
the semi-diagonal cubic sector requires a number of displaced-gradient
evaluations growing only linearly with the number of vibrational modes,
whereas the genuinely three-index sector requires a number of evaluations
growing quadratically. The split of the sextic cubic contribution into
semi-diagonal and three-index parts therefore reflects a real difference in
computational scaling, not just a formal reorganization of the equations.\cite{cubic-vib_2024,fast_dvib_2025}

This elimination of primitive $d^2I/dQ^2$ is essential for the main text. It
is precisely what allows standard quartics and standard sextics to belong to
the same linear representation-equivalent sector.

\section{Intrinsic Second-Level Quartic Tensor and the \texorpdfstring{$H_{22}$}{H22} Channel}

\subsection{Decomposition of the second inverse-inertia derivative}

The first genuine post-standard quartic object enters through the second
inverse-inertia derivative,
\begin{equation}
\mu_{2,kl} \equiv \frac{\partial^2 I^{-1}}{\partial Q_k \partial Q_l}.
\end{equation}
This quantity admits a natural decomposition into an algebraic bilinear closure
of the first-level tensor and a genuinely new intrinsic remainder.

From the identity
\begin{equation}
\frac{\partial I^{-1}}{\partial Q_k}
=
-I^{-1}\frac{\partial I}{\partial Q_k}I^{-1},
\end{equation}
a second differentiation yields
\begin{align}
\mu_{2,kl}
&=
I^{-1}\frac{\partial I}{\partial Q_k}I^{-1}\frac{\partial I}{\partial Q_l}I^{-1}
+
I^{-1}\frac{\partial I}{\partial Q_l}I^{-1}\frac{\partial I}{\partial Q_k}I^{-1}
-
I^{-1}\frac{\partial^2 I}{\partial Q_k\partial Q_l}I^{-1}.
\end{align}
Writing
\begin{equation}
\mu_{1,k}=\frac{\partial I^{-1}}{\partial Q_k},
\end{equation}
this becomes
\begin{equation}
\mu_{2,kl} = B_{kl}(\mu_1) - N_{kl},
\label{eq:app_mu2_split}
\end{equation}
with
\begin{equation}
B_{kl}(\mu_1)=\mu_{1,k}I\mu_{1,l}+\mu_{1,l}I\mu_{1,k}
\end{equation}
and
\begin{equation}
N_{kl}=I^{-1}\frac{\partial^2 I}{\partial Q_k\partial Q_l}I^{-1}.
\end{equation}

Equation~\eqref{eq:app_mu2_split} is one of the central structural identities of
the present work. It shows that the full second inverse-inertia derivative
contains both:
\begin{enumerate}
\item an algebraic closure already generated by the first-level tensor $\mu_1$,
\item a genuinely new intrinsic second-level object $N$.
\end{enumerate}

\subsection{Quartic channel \texorpdfstring{$H_{22}$}{H22}}

The quartic channel $H_{22}$ is the first contribution for which the new
intrinsic tensor $N$ is unavoidable. In the implementation, the full $H_{22}$
contribution is written as a weighted quadratic form in the symmetric mode-pair
tensor $\mu_2$, or equivalently as
\begin{equation}
H_{22} = H_{22}[B(\mu_1)-N,\;B(\mu_1)-N].
\end{equation}
Expanding this expression gives the operational decomposition
\begin{equation}
H_{22}
=
H_{22}[B,B]
-H_{22}[B,N]
-H_{22}[N,B]
+H_{22}[N,N].
\end{equation}
This is the explicit sense in which $H_{22}$ is the first quartic breakdown
term. Its bilinear part still belongs to the closure of the standard
first-level sector, whereas the intrinsic part is controlled by $N$ and
therefore lies outside the standard linear framework.

\subsection{Covariance of the intrinsic tensor}

The intrinsic second-level tensor $N$ transforms covariantly under axis
permutation, exactly as required for a genuine tensorial object. In the code,
this covariance was checked numerically by rebuilding the same harmonic model in
different principal-axis representations and verifying agreement of
\begin{equation}
\mu_1,\qquad B(\mu_1),\qquad N,\qquad \mu_2
\end{equation}
up to the expected representation permutations and mode-sign conventions.

This point matters because it shows that the breakdown introduced by $H_{22}$
is not a numerical artifact of a particular coordinate choice. The new object
$N$ is a bona fide tensorial quantity with a well-defined transformation law.

\section{First Sextic Post-Standard Diagnostic: Linear Response in \texorpdfstring{$\tau(H_{22})$}{tau(H22)}}

\subsection{Why the sextic analogue should be sought in the \texorpdfstring{$\tau$}{tau} channel}

The standard sextic geometry functional depends explicitly on the quartic
Wilson tensor. Since the standard quartic part belongs to the first-level
sector, the simplest post-standard sextic effect is not expected to arise from
a new primitive sextic block of the same simplicity as $H_{22}$. Rather, the
first natural candidate is obtained by inserting the non-standard quartic
correction generated by $H_{22}$ into the standard sextic geometry functional.

Accordingly, one writes
\begin{equation}
\tau = \tau_{\mathrm{std}} + \delta\tau,
\qquad
\delta\tau \equiv \tau(H_{22}),
\end{equation}
and defines the associated sextic response by isolating the term linear in
$\delta\tau$.

\subsection{Linear response extraction}

If $G_\tau$ denotes the part of the sextic geometry functional depending on the
quartic tensor, the first post-standard sextic diagnostic is
\begin{equation}
\Delta^{(6)}_{H_{22}}
=
G_\tau(\tau_{\mathrm{std}}+\delta\tau)
-G_\tau(\tau_{\mathrm{std}})
-G_\tau(\delta\tau).
\label{eq:app_sextic_linear_response}
\end{equation}
Equation~\eqref{eq:app_sextic_linear_response} is the symmetric bilinear cross
term between the standard quartic sector and the quartic correction generated
by $H_{22}$. In the code this quantity is evaluated component by component for
the full set of sextic Cartesian coefficients.

\subsection{Interpretation}

The quantity $\Delta^{(6)}_{H_{22}}$ is not intended here as a complete
post-standard sextic correction. Its role is diagnostic. It is:
\begin{enumerate}
\item the first sextic contribution that feels the new quartic tensorial object
generated by $H_{22}$;
\item computable from the same basic quantities already required in the
standard machinery;
\item simple enough to be used as a practical indicator of the onset of the
breakdown of the standard linear sector.
\end{enumerate}

For H$_2$O, the largest component of this sextic diagnostic lies at the level
of about $10^3$ Hz. This is small compared with the dominant standard sextic
constants, but large enough to define a clear post-standard scale. The same
construction is expected to become smaller in systems for which standard
VPT-based and higher-level treatments agree more closely.

\section{Computational Diagnostics and Practical Use}

The structure derived above has immediate practical consequences. Because the
quartic diagnostic $H_{22}$ and the sextic diagnostic
$\Delta^{(6)}_{H_{22}}$ are built from quantities already present in the
standard formalism, they can be added to tensor-based analysis tools without
requiring a qualitatively new electronic-structure workflow.

In practical terms, this means that:
\begin{enumerate}
\item standard quartic and sextic constants may be transported between
representations by the pseudoinverse-lift/tensor-transport/reprojection
procedure;
\item the magnitude of $H_{22}$ may be monitored as the first quartic indicator
that the standard linear sector is beginning to fail;
\item the magnitude of $\Delta^{(6)}_{H_{22}}$ may be monitored as the first
sextic analogue of the same phenomenon.
\end{enumerate}
These diagnostics are therefore natural candidates to complement traditional
quantities such as $s_{111}$ in practical fitting strategies and in the app
associated with the present framework.
